\documentclass[11pt,a4paper]{report}
% Préambule partagé des rapports SecuriGlass Factory OS.
% Style volontairement sobre : pas d'encadrés de couleur, pas de blocs
% décoratifs. Règles horizontales fines, typographie classique, une seule
% couleur d'accent reprise de l'identité visuelle du produit.

\usepackage{fontspec}
\usepackage{polyglossia}
\setdefaultlanguage{french}
\setmainfont{Palatino}
\setsansfont{Helvetica Neue}
\setmonofont{Menlo}[Scale=0.85]
\usepackage{microtype}
\sloppy
\usepackage[a4paper,top=2.8cm,bottom=2.8cm,left=2.6cm,right=2.6cm,headsep=1.1cm]{geometry}
\usepackage{graphicx}
\usepackage{booktabs}
\usepackage{longtable}
\usepackage{array}
\usepackage{caption}
\usepackage{enumitem}
\usepackage{xcolor}
\usepackage{fancyhdr}
\usepackage{titlesec}
\usepackage{listings}
\usepackage{tocloft}
\usepackage[hidelinks]{hyperref}
\usepackage[nameinlink,french]{cleveref}

% ---------------------------------------------------------------- couleurs
\definecolor{accent}{HTML}{15584B}
\definecolor{accentlight}{HTML}{1F7A68}
\definecolor{graytext}{HTML}{3F4A47}
\definecolor{rulegray}{HTML}{C9CFCD}

\hypersetup{
  colorlinks=true,
  linkcolor=accent,
  citecolor=accent,
  urlcolor=accentlight,
  pdfauthor={Mohamed Aziz Baoueb},
}

% ------------------------------------------------------------------ titres
\titleformat{\chapter}[hang]
  {\normalfont\Huge\bfseries\color{accent}}
  {\thechapter}{1.2cm}{}
\titlespacing*{\chapter}{0pt}{-0.8cm}{1.0cm}

\titleformat{\section}
  {\normalfont\Large\bfseries\color{accent}}
  {\thesection}{0.6em}{}
\titleformat{\subsection}
  {\normalfont\large\bfseries\color{graytext}}
  {\thesubsection}{0.6em}{}
\titleformat{\subsubsection}
  {\normalfont\normalsize\bfseries\itshape\color{graytext}}
  {\thesubsubsection}{0.6em}{}

% --------------------------------------------------------------- en-têtes
\pagestyle{fancy}
\fancyhf{}
\renewcommand{\headrulewidth}{0.4pt}
\renewcommand{\footrulewidth}{0pt}
\fancyhead[L]{\small\itshape\nouppercase{\leftmark}}
\fancyhead[R]{\small\thepage}
\fancypagestyle{plain}{
  \fancyhf{}
  \fancyhead[R]{\small\thepage}
  \renewcommand{\headrulewidth}{0.4pt}
}

% ----------------------------------------------------------------- listes
\setlist{itemsep=2pt, topsep=4pt, parsep=0pt}

% --------------------------------------------------------------- légendes
\captionsetup{
  labelfont={bf,color=accent},
  font=small,
  skip=6pt,
  justification=centering,
}

% ----------------------------------------------------------- code source
\lstset{
  basicstyle=\ttfamily\small,
  breaklines=true,
  frame=none,
  columns=flexible,
  keywordstyle=\color{accent}\bfseries,
  commentstyle=\color{graytext}\itshape,
  stringstyle=\color{accentlight},
  showstringspaces=false,
  tabsize=2,
}

% -------------------------------------------------------- table des matières
\renewcommand{\cftchapfont}{\bfseries\color{accent}}
\renewcommand{\cftchappagefont}{\color{accent}}

% ----------------------------------------------------------- environnement
% figure pleine largeur avec légende, sans cadre ni fond.
\newcommand{\screenshot}[3]{%
  \begin{figure}[htbp]
    \centering
    \includegraphics[width=#2\textwidth]{figures/#1}
    \caption{#3}
  \end{figure}%
}

% séparateur fin entre deux sections d'un même thème (jamais de remplissage).
\newcommand{\finesep}{%
  \par\noindent\textcolor{rulegray}{\rule{\linewidth}{0.4pt}}\par\vspace{4pt}%
}


\title{SecuriGlass Factory OS --- Rapport technique}
\author{Mohamed Aziz Baoueb}
\date{Octobre 2026}

\begin{document}

% ============================================================ page de titre
\begin{titlepage}
  \centering
  \vspace*{2.2cm}

  {\Large\scshape Rapport technique complet}\\[0.6cm]
  \textcolor{rulegray}{\rule{0.5\linewidth}{0.4pt}}\\[1.4cm]

  {\fontsize{32}{38}\selectfont\bfseries\color{accent} SecuriGlass Factory~OS}\\[0.5cm]
  {\Large\itshape Architecture, implémentation, sécurité et vérification}\\[2.2cm]

  \begin{tabular}{rl}
    \textbf{Auteur} & Mohamed Aziz Baoueb \\[4pt]
    \textbf{Date} & Octobre 2026 \\[4pt]
    \textbf{Nature} & Documentation technique de preuve de concept (PoC) \\[4pt]
    \textbf{Langages} & Python, TypeScript, C, SQL, Bash \\[4pt]
    \textbf{Référentiels} & IEC~62443-3-3, OWASP ASVS v4.0, OWASP Top~10 for LLM,\\
     & NIST SP~800-82, RGPD, EN~12150-1, ANSSI EBIOS~RM \\
  \end{tabular}

  \vfill
  \textcolor{rulegray}{\rule{0.5\linewidth}{0.4pt}}\\[0.4cm]
  {\small Ce document détaille l'implémentation technique du projet~: pile
  logicielle, architecture, sécurité, tests et déploiement. Le rapport de
  projet, plus court, présente les objectifs et les fonctionnalités.}
\end{titlepage}

\tableofcontents

% ================================================================ chapitre
\chapter{Introduction}

\section{Objet de ce document}
Ce rapport documente, de façon exhaustive, l'implémentation technique de
SecuriGlass Factory~OS~: la pile logicielle retenue pour chaque composant, la
structure du dépôt, les mécanismes de sécurité réellement codés (et non
seulement envisagés), la stratégie de tests, et la manière de faire tourner
l'ensemble en local. Il complète le rapport de projet, qui reste centré sur les
objectifs et les fonctionnalités vues par l'utilisateur.

\section{Principe de rédaction}
Chaque mesure de sécurité décrite ici a été vérifiée dans le code au moment de
la rédaction. Quand une mesure n'est pas implémentée mais seulement prévue,
elle est explicitement marquée \emph{planifiée} ; ce choix évite qu'une lecture
rapide du document fasse croire à une garantie qui n'existe pas encore dans le
code.

\chapter{Architecture générale}

\section{Vue d'ensemble}
Le système est composé de cinq services indépendants, chacun dans son propre
répertoire sous \texttt{services/}, communiquant par des interfaces strictement
définies (Modbus~TCP, UDP, HTTP, WebSocket). Aucun composant ne partage de
mémoire ou de base de données avec un autre composant d'une zone de confiance
différente.

\begin{center}
\begin{tabular}{@{}lll@{}}
\toprule
\textbf{Service} & \textbf{Rôle} & \textbf{Langage} \\
\midrule
\texttt{ot-plc-daemon} & Automate du four (simulateur + serveur Modbus) & C11 \\
\texttt{ot-edge-gateway} & Passerelle de collecte OT~$\rightarrow$~diode & Python 3.12 \\
\texttt{backend-api} & API métier, authentification, IA, WebSocket & Python 3.12 / FastAPI \\
\texttt{frontend-ui} & Console web & TypeScript / React 18 \\
\texttt{soc-siem} & Détection et corrélation (hors périmètre console) & Python \\
\bottomrule
\end{tabular}
\end{center}

\section{Le modèle de Purdue appliqué}
Chaque service est rattaché à un niveau du modèle de Purdue, qui organise un
système industriel en zones de confiance croissante~:

\begin{center}
\begin{tabular}{@{}cl p{7.5cm}@{}}
\toprule
\textbf{Niveau} & \textbf{Zone} & \textbf{Composants et contraintes} \\
\midrule
0--1 & Terrain / contrôle & \texttt{ot-plc-daemon} : aucune dépendance réseau
sortante, aucune écriture distante non authentifiée. \\
2--3 & Supervision OT & \texttt{ot-edge-gateway} : client Modbus en
\emph{lecture seule} uniquement (aucune fonction d'écriture implémentée côté
client). \\
3,5 & Zone démilitarisée & Diode de données : liaison UDP strictement
unidirectionnelle, aucun accusé de réception possible par construction du
protocole. \\
4 & Informatique de gestion & \texttt{backend-api} + PostgreSQL : seul point
d'entrée du système d'information vers les données OT. \\
5 & Entreprise & \texttt{frontend-ui} : aucun accès réseau direct à l'OT, tout
passe par l'API de niveau~4. \\
\bottomrule
\end{tabular}
\end{center}

La propriété de sécurité centrale du système est que \textbf{le franchissement
du niveau 3,5 vers le niveau~3 est impossible par construction}~: la diode de
données est un émetteur UDP \emph{sans aucune fonction de réception}, donc même
une compromission complète du niveau~4 ne donnerait pas, à elle seule, de
chemin réseau vers l'automate.

\chapter{Pile technologique}
\label{chap:tech-stack}

\section{Backend}
\begin{itemize}
  \item \textbf{FastAPI} (Python 3.12, async) --- framework HTTP et WebSocket.
  \item \textbf{asyncpg} --- pilote PostgreSQL asynchrone, sans ORM, requêtes
    SQL paramétrées explicites.
  \item \textbf{Pydantic~v2} --- validation stricte des schémas d'entrée/sortie
    (\texttt{extra="forbid"}, rejet de tout champ inconnu).
  \item \textbf{Argon2id} (bibliothèque \texttt{argon2-cffi}) --- hachage des
    mots de passe, paramètres $t=3$, $m=65536$~KiB (64~Mio), $p=4$, sel de
    16~octets, conformes à la RFC~9106.
  \item \textbf{PyJWT} --- jetons de session signés HS256, identifiant unique
    (\texttt{jti}) permettant la révocation.
  \item \textbf{cryptography} (AES-256-GCM) --- chiffrement au repos des
    données personnelles clients (TDE applicative).
  \item \textbf{scikit-learn} (repli pur Python si absent) --- détection
    d'anomalies par forêt d'isolement.
  \item \textbf{Alembic} --- migrations de schéma versionnées et idempotentes.
  \item \textbf{mypy --strict}, \textbf{bandit}, \textbf{black} --- qualité et
    sécurité statique du code.
\end{itemize}

\section{Frontend}
\begin{itemize}
  \item \textbf{React~18} + \textbf{TypeScript} (mode strict,
    \texttt{noUncheckedIndexedAccess}, \texttt{exactOptionalPropertyTypes}).
  \item \textbf{Vite} --- empaquetage ; build de production sans script
    inline, compatible avec une Content-Security-Policy stricte.
  \item \textbf{Zustand} --- état applicatif, sans middleware de persistance
    (le jeton de session ne vit qu'en mémoire JavaScript).
  \item \textbf{react-router} --- routage côté client, gardes d'accès par
    rôle.
  \item \textbf{Vitest} + \textbf{Testing Library} --- tests unitaires et de
    composants, couverture~\textgreater{}~80~\% imposée en seuil de build.
  \item Aucune bibliothèque de rendu HTML dynamique, aucune dépendance
    graphique tierce~: les graphiques sont du SVG généré par le code.
\end{itemize}

\section{Automate et réseau industriel}
\begin{itemize}
  \item \textbf{C11} (\texttt{-std=c11 -Wall -Wextra -Werror -pedantic}) ---
    démon automate, zéro allocation dynamique résiduelle, \texttt{poll()} pour
    la boucle réseau non bloquante.
  \item \textbf{Modbus Application Protocol V1.1b3} --- codes fonction~0x01 à
    0x10 implémentés côté serveur (lecture de bobines, de registres, écriture
    simple et multiple).
  \item \textbf{Python 3.12} (passerelle) --- client Modbus en lecture seule,
    \texttt{cryptography} pour la signature Ed25519 des trames.
\end{itemize}

\section{Infrastructure}
\begin{itemize}
  \item \textbf{Docker} --- conteneurisation de PostgreSQL (avec extension
    \texttt{pgAudit}) et des services applicatifs.
  \item \textbf{Ansible} --- durcissement système conforme aux CIS Benchmarks
    (rôle \texttt{cis\_linux}).
  \item \textbf{Terraform} --- définition déclarative de l'infrastructure
    cible.
  \item \textbf{GitHub Actions} --- intégration continue en sept étapes
    (\cref{chap:qualite}).
\end{itemize}

\chapter{Structure du dépôt}

\begin{lstlisting}[language={}]
poc-glass/
  database/                 Schemas SQL, migrations Alembic
    schemas/                01_init .. 06_factory_extension.sql
    migrations/versions/    0001 .. 0004
  docs/                     Architecture, GRC, specifications, QA
  infra/
    docker/                 Dockerfile PostgreSQL + docker-compose
    ansible/                Durcissement CIS (rôle cis_linux)
    terraform/               Infrastructure déclarative
  redteam/                  Outils de test d'intrusion Modbus (lecture
                             et écriture encadrées, autorisation requise)
  scripts/                  dev_up.sh / dev_down.sh / seeds de démo
  security/                 Config pipeline SAST/SCA, tests associés
  services/
    backend-api/            API FastAPI (app/api, app/core, app/services)
    frontend-ui/             Console React/TypeScript (src/)
    ot-plc-daemon/           Automate C11 (src/, include/, tests/)
    ot-edge-gateway/         Passerelle Python (gateway/)
    soc-siem/                Détection et corrélation
  tests/                    Tests d'intégration et de sécurité transverses
\end{lstlisting}

\chapter{Backend --- API et services métier}

\section{Organisation du code}
Le backend suit une séparation stricte en trois couches~: \texttt{app/api}
(routes HTTP/WebSocket, validation d'entrée, dépendances FastAPI),
\texttt{app/services} (logique métier pure, testable indépendamment du
framework web) et \texttt{app/core} (configuration, base de données,
sécurité transverse).

\section{Identité et authentification}
Les mots de passe sont hachés en Argon2id (paramètres au
\cref{chap:tech-stack}). La connexion délivre un jeton JWT (HS256) contenant un
identifiant unique de session (\texttt{jti}) ; ce \texttt{jti} est inscrit dans
une table de révocation lors de la déconnexion, ce qui permet d'invalider un
jeton avant son expiration naturelle sans avoir besoin d'une liste noire
globale à chaque requête. Un verrouillage temporaire du compte suit plusieurs
échecs de connexion consécutifs.

\section{Autorisation~: RBAC puis ABAC}
Six rôles sont définis~: \texttt{SUPER\_ADMIN}, \texttt{SECURITY\_AUDITOR},
\texttt{PLANT\_MANAGER}, \texttt{OT\_OPERATOR}, \texttt{QUALITY\_INSPECTOR} et
\texttt{SERVICE\_ACCOUNT\_EDGE}. Chaque route de l'API déclare explicitement
les rôles autorisés via une dépendance FastAPI (\texttt{require\_roles}), et
toute tentative refusée est inscrite dans le journal d'audit avec le rôle
détenu et le rôle requis. Un moteur de règles contextuelles (ABAC) vient
compléter le RBAC pour des décisions qui dépendent du contexte (plage horaire,
sous-réseau, verrou de maintenance) : la politique est \emph{fail-closed}~:
en l'absence de donnée de contexte, la règle refuse plutôt que d'autoriser par
défaut.

\section{Confidentialité des données clients}
Les champs personnels des clients (nom du contact, e-mail, téléphone) sont
chiffrés avec AES-256-GCM avant écriture en base ; seule la raison sociale
reste en clair, pour permettre la recherche sans exposer de donnée
personnelle. Le déchiffrement n'a lieu qu'à la lecture par un rôle habilité.

\section{Journal d'audit chaîné}
Chaque événement sensible (connexion, création de client, refus d'accès,
consigne bloquée, vérification de la chaîne elle-même) est inscrit dans
\texttt{sec\_data.audit\_logs} avec une empreinte SHA-256 calculée sur
l'enregistrement \emph{et} l'empreinte de l'enregistrement précédent. Un
verrou consultatif PostgreSQL (\texttt{pg\_advisory\_xact\_lock}) sérialise les
écritures concurrentes pour que la chaîne reste strictement ordonnée. Un
déclencheur SQL interdit toute modification ou suppression de ligne existante
\texttt{a posteriori} ; la suppression d'un compte (exigée par le RGPD,
art.~17) ne touche donc jamais le journal, qui conserve un instantané de
l'identité de l'acteur sans contrainte référentielle vers la table des
utilisateurs.

\section{Télémétrie en direct~: WebSocket}
Le canal \texttt{/ws/telemetry} authentifie le client dans le \emph{premier
message} après l'ouverture de la connexion (jamais dans l'URL, pour éviter
qu'un jeton se retrouve dans des journaux d'accès). Le rôle est revérifié
périodiquement (toutes les 30~secondes) pendant la durée de la connexion, et
le nombre de connexions simultanées est plafonné. Chaque nouvelle mesure est
accompagnée du diagnostic du modèle d'apprentissage automatique, jamais
présentée comme « normale » par défaut quand ce diagnostic est
indisponible.

\section{Détection d'anomalies (apprentissage automatique)}
Le détecteur utilise une forêt d'isolement (\emph{Isolation Forest},
\texttt{scikit-learn}, avec un repli pur Python si la bibliothèque est
absente) sur neuf mesures physiques (températures des quatre zones,
pyromètre du verre, pression de trempe, vitesse de convoyeur, courant du
soufflage, puissance consommée) enrichies de trois grandeurs dérivées~: le
régime de fonctionnement courant, l'écart maximal entre zones, et l'écart
entre le pyromètre et la moyenne des zones. Ces grandeurs dérivées donnent au
modèle le contexte croisé qu'une lecture zone par zone ne porte pas. En
complément du score statistique, des signatures de pannes déterministes
(surchauffe critique, perte de pression de trempe, blocage convoyeur, risque
de casse thermique, défaillance d'une résistance) forcent une sévérité
\texttt{CRITICAL} indépendamment du score, pour garantir leur détection à
100~\% quelle que soit la sensibilité statistique du modèle.

\section{Assistant technique et garde LLM}
\label{sec:llm-guard}
L'assistant est protégé par quatre barrières hermétiques, appliquées dans
l'ordre~:
\begin{enumerate}
  \item \textbf{Filtrage de l'entrée} --- canonicalisation du texte (décodage
    des échappements HTML/URL, normalisation Unicode NFKC, suppression des
    caractères de contrôle, mise à plat des variantes à base d'espacement ou
    de \emph{leetspeak}, décodage récursif borné du Base64) avant
    comparaison à un ensemble de motifs d'attaque multilingues (français,
    anglais, espagnol)~: contournement d'instruction, désactivation de
    sécurité, jeu de rôle, exfiltration de secret ou de message système.
  \item \textbf{Recherche cloisonnée} --- les réponses sont puisées dans un
    corpus documentaire local portant sur la norme EN~12150, jamais générées
    librement à partir de connaissances externes.
  \item \textbf{Génération contrainte} --- le moteur de réponse est
    déterministe et contraint au contexte normatif récupéré ; il n'improvise
    pas de valeur physique.
  \item \textbf{Contrôle de sortie} --- toute grandeur physique présente dans
    la réponse (température, pression, quelle que soit l'unité employée) est
    reconvertie et comparée aux bornes de sûreté EN~12150
    (620--730\textdegree{}C, 1,5--6,5~bar) avant envoi ; une réponse hors
    bornes est bloquée et consignée comme incident critique.
\end{enumerate}
Un limiteur de débit à seau de jetons (10~requêtes/minute/utilisateur)
protège l'assistant contre un abus par volume. Les demandes portant sur des
données financières ou personnelles des clients sont refusées par un filtre
d'intention dédié, appliqué \emph{après} le filtre anti-injection, et
consignées séparément.

\section{API « usine » et principe du besoin d'en connaître}
L'endpoint \texttt{GET~/api/v1/factory/overview} ne renvoie que les blocs de
données utiles au rôle de l'appelant~: un opérateur ne reçoit ni les données
commerciales, ni la conformité, ni les alertes, même si le serveur les a
calculées pour un autre appel. Ce filtrage est appliqué côté serveur, pas
seulement masqué côté interface. La prise de commande
(\texttt{POST~/api/v1/orders}) exécute, dans une seule transaction~: la
vérification du stock disponible (verrouillage des lignes concernées), le
calcul de la surface de verre, la déduction des matières premières et
l'enregistrement de la commande ; tout échec annule l'ensemble.

\chapter{Réseau industriel~: automate et passerelle}

\section{Simulateur physique et serveur Modbus}
Le démon \texttt{ot-plc-daemon} simule la thermodynamique d'un four à quatre
zones de chauffe (inertie différentielle, refroidissement par soufflage,
cinématique du convoyeur) et expose ces grandeurs via un serveur Modbus~TCP
non bloquant (\texttt{poll()}, un seul processus, aucune allocation
dynamique résiduelle). Les consignes acceptées en écriture sont bornées
physiquement dans le micrologiciel lui-même (ex.~500,0--750,0\textdegree{}C
pour les consignes de zone) ; une valeur hors bornes est refusée par une
exception Modbus (\texttt{ILLEGAL\_DATA\_VALUE}), indépendamment de tout
contrôle applicatif en amont. L'arrêt d'urgence (E-STOP) coupe le chauffage
immédiatement ; sa levée exige un acquittement physique du relais de
sécurité et ne peut pas être demandée par le réseau seul, par analogie avec
un automate réel.

\section{Passerelle de collecte}
La passerelle \texttt{ot-edge-gateway} interroge l'automate en Modbus avec un
client strictement \emph{read-only} (aucune fonction d'écriture n'est même
implémentée côté client), décode les registres en grandeurs physiques, les
valide (bornes physiques, absence de valeurs non finies), puis les transmet
à la diode de données. Elle ne reçoit jamais de réponse du système
d'information.

\section{Diode de données}
Le protocole binaire de la diode comprend un en-tête fixe de 64~octets
(signature \texttt{SGDD}, version, numéro de séquence croissant sur 64~bits,
horodatage en microsecondes, identifiant de l'émetteur, longueur de charge
utile), suivi d'une charge utile JSON canonique et d'un sceau de 96~octets
(empreinte SHA-256 et signature Ed25519). Le récepteur, côté système
d'information~:
\begin{itemize}
  \item recalcule l'empreinte en temps constant (\texttt{hmac.compare\_digest})
    et vérifie la signature avec la clé publique de confiance de l'émetteur~;
  \item rejette toute trame dont la séquence n'est pas strictement croissante
    (protection anti-rejeu) et tolère une dérive d'horloge maximale de
    30~secondes~;
  \item déduplique durablement par un identifiant de transmission unique
    (UUID~v5) contraint \texttt{UNIQUE} en base, pour rejeter un rejeu même
    après un redémarrage de la passerelle~;
  \item journalise toute trame rejetée dans
    \texttt{ot\_data.diode\_import\_log}, sans jamais émettre d'accusé de
    réception vers le réseau industriel.
\end{itemize}

\chapter{Interface web}

\section{Sécurité côté client}
La politique de sécurité du contenu (CSP) interdit tout script non
auto-hébergé et tout script en ligne (\texttt{default-src~'none'}) ; le
jeton de session ne vit qu'en mémoire JavaScript (aucune écriture dans
\texttt{localStorage}, \texttt{sessionStorage} ou un cookie), ce qui est
vérifié automatiquement par un test qui interdit la présence de ces appels
dans l'ensemble du code source. Un rechargement de page déconnecte donc
l'utilisateur par construction.

\section{Accès par rôle, dans l'interface comme sur le réseau}
La navigation, les routes et les sélecteurs d'état n'affichent une rubrique
que si le rôle courant y donne droit ; ce filtrage côté client est une
commodité d'usage, jamais une autorité~: chaque appel réseau correspondant
est revérifié côté serveur, et une tentative d'accès direct à une route non
autorisée renvoie une page d'erreur sans qu'aucune donnée n'ait transité.

\section{Thème et langue}
L'interface suit la préférence système pour le mode clair ou sombre
(\texttt{prefers-color-scheme}) et propose une bascule manuelle ; la console
est disponible en français et en anglais via un dictionnaire de traduction
central. Ces deux préférences restent en mémoire pour la durée de la session,
sans être écrites sur le poste, pour respecter la même contrainte que le
jeton de session.

\section{Accessibilité et robustesse d'affichage}
Les composants graphiques (courbes, jauges, carte réseau) sont du SVG
généré par le code, jamais une image~: aucune dépendance à un style en ligne
interdit par la CSP. La mise en page reste lisible depuis un téléphone
(375~px) jusqu'à un grand écran, avec un empilement explicite des éléments
plutôt qu'une grille figée susceptible de faire chevaucher du texte long
(ex.~nom de client) sur une colonne voisine.

\chapter{Qualité, tests et intégration continue}
\label{chap:qualite}

\section{Stratégie de tests}
\begin{center}
\begin{tabular}{@{}llr@{}}
\toprule
\textbf{Domaine} & \textbf{Outil} & \textbf{Résultat} \\
\midrule
API et services métier & pytest (avec PostgreSQL réel) & 563 passés \\
Interface web & Vitest + Testing Library & 122 passés \\
Automate C --- unitaires & \texttt{bin/plc\_test} & 22 passés \\
Automate C --- fuzzing & \texttt{bin/plc\_fuzz\_test} & 763 cas, 0 crash \\
Automate C --- fuites mémoire & \texttt{leaks} (macOS) & 0 fuite \\
Passerelle de collecte & pytest & 20 passés \\
Red team diode / LLM Guard / ML & pytest (adversarial) & 133 passés \\
\bottomrule
\end{tabular}
\end{center}

\section{Portes de qualité automatisées}
L'intégration continue (GitHub Actions) exécute sept étapes à chaque
modification~: \texttt{lint} (types stricts, formatage), \texttt{test}
(suite complète), \texttt{sast} (analyse statique de sécurité, \texttt{bandit}
et règles maison), \texttt{sca} (audit des dépendances), \texttt{docker}
(construction des images), \texttt{sbom} (inventaire logiciel) et
\texttt{sign} (signature des artefacts produits).

\section{Vérification par attaque réelle}
Au-delà des tests automatisés, le comportement du système a été vérifié par
une attaque réseau réelle menée directement contre l'automate (sans passer
par l'application), décrite au chapitre~6 du rapport de projet~: elle a
confirmé que les bornes physiques tiennent au niveau du micrologiciel, que le
modèle de détection réagit sans intervention, et qu'elle a aussi révélé
l'absence de sonde réseau sur le segment industriel --- une limite
honnêtement documentée plutôt que masquée.

\chapter{Infrastructure et déploiement}

\section{Environnement de développement local}
Le script \texttt{scripts/dev\_up.sh} orchestre le démarrage complet de la
chaîne en sept étapes~: base PostgreSQL durcie (avec \texttt{pgAudit}),
génération d'une paire de clés Ed25519 pour la passerelle, démarrage de
l'API, création des comptes et données de démonstration, compilation et
démarrage de l'automate, démarrage du récepteur et de la passerelle, puis
démarrage de la console. \texttt{scripts/dev\_down.sh} arrête proprement
l'ensemble.

\section{Durcissement système}
Le rôle Ansible \texttt{cis\_linux} applique un sous-ensemble des CIS
Benchmarks~: durcissement \texttt{sysctl} (désactivation du
\emph{forwarding} IP, protection anti-usurpation, activation des
\emph{SYN cookies}, désactivation des redirections ICMP), configuration de
\texttt{sshd} et d'\texttt{auditd}, bannière légale de connexion.

\section{Conteneurisation}
PostgreSQL tourne dans une image Debian (et non Alpine, qui ne fournit pas
l'extension \texttt{pgAudit}) avec principe de moindre privilège
(\texttt{no-new-privileges}, capacités Linux réduites au strict nécessaire).

\chapter{Limites connues et travaux futurs}

\section{Ce qui est implémenté}
Tout ce qui est décrit dans ce rapport correspond à du code existant et
testé au moment de la rédaction.

\section{Ce qui est volontairement hors périmètre}
\begin{itemize}
  \item \textbf{Détection réseau sur le segment OT} --- une attaque Modbus
    directe sur l'automate n'est pas aujourd'hui détectée par le flux
    d'alertes de la console (voir chapitre~6 du rapport de projet).
  \item \textbf{Réplication à froid et plan de reprise} --- cités dans le
    dossier d'architecture comme mesures cibles, non implémentés dans ce
    PoC.
  \item \textbf{Signature des consignes envoyées à l'automate} --- l'écriture
    de consignes existe au niveau du protocole Modbus pour les besoins de la
    démonstration, sans couche de signature applicative supplémentaire.
  \item \textbf{Image Docker \emph{distroless}} pour le backend --- non
    construite dans ce PoC.
  \item \textbf{Homologation} --- ce projet est une démonstration technique ;
    aucune autorité d'homologation ne s'est prononcée sur le système.
\end{itemize}

\section{Suites possibles}
Ajouter une sonde de détection sur le réseau industriel, exécuter
l'intégralité de la suite de tests (y compris ceux aujourd'hui ignorés faute
d'outillage disponible sur macOS) sur un environnement Linux de référence,
puis traiter dans l'ordre les mesures listées ci-dessus comme
\emph{planifiées}.

\appendix
\chapter{Glossaire}

\begin{description}
  \item[ABAC] \emph{Attribute-Based Access Control} --- autorisation fondée
    sur des règles contextuelles, en complément du contrôle par rôle.
  \item[CSP] \emph{Content Security Policy} --- politique de sécurité du
    contenu, appliquée par le navigateur.
  \item[Diode de données] Dispositif garantissant qu'un flux d'information ne
    peut transiter que dans un seul sens entre deux zones réseau.
  \item[IEC~62443] Famille de normes de cybersécurité pour les systèmes
    d'automatisation et de contrôle industriels.
  \item[Modèle de Purdue] Modèle de référence organisant un système
    industriel en niveaux de confiance, du terrain (niveau~0) à l'entreprise
    (niveau~5).
  \item[Modbus~TCP] Protocole de communication industrielle utilisé pour
    interroger et commander un automate.
  \item[OT / IT] \emph{Operational Technology} (réseau industriel) et
    \emph{Information Technology} (réseau informatique de gestion).
  \item[RBAC] \emph{Role-Based Access Control} --- autorisation fondée sur le
    rôle de l'utilisateur.
  \item[TDE] \emph{Transparent Data Encryption} --- chiffrement des données
    au repos, transparent pour l'application qui les consomme.
\end{description}

\end{document}
